\documentclass[11pt,a4paper]{article}
\usepackage[utf8]{inputenc}
\usepackage[T1]{fontenc}
\usepackage[french]{babel}
\usepackage{geometry}
\geometry{margin=2.5cm}
\usepackage{graphicx}
\usepackage{xcolor}
\usepackage{amsmath}
\usepackage{listings}
\usepackage{hyperref}
\usepackage{fancyhdr}
\usepackage{titlesec}
\usepackage{enumitem}
\usepackage{tcolorbox}
\usepackage{float}
\usepackage{tikz}
\usepackage{tabularx}
\usepackage{longtable}
\usepackage{booktabs}
\usetikzlibrary{shapes.geometric,arrows.meta,positioning,calc,fit,backgrounds,decorations.pathreplacing}

% Configuration des listes
\setlist[itemize]{leftmargin=*, itemsep=4pt, topsep=4pt}
\setlist[enumerate]{leftmargin=*, itemsep=4pt, topsep=4pt}

% Configuration du code
\lstset{
    basicstyle=\ttfamily\small,
    breaklines=true,
    frame=single,
    keywordstyle=\color{blue},
    commentstyle=\color{green!60!black},
    stringstyle=\color{orange},
    tabsize=4,
    showstringspaces=false
}

% En-tête et pied de page
\pagestyle{fancy}
\fancyhf{}
\fancyhead[L]{\small Documentation Phase 2 --- ChatBot SUP'ONE}
\fancyhead[R]{\small \today}
\fancyfoot[C]{\thepage}

% Environnements personnalisés
\newtcolorbox{important}[1][]{
  colback=red!5!white,
  colframe=red!75!black,
  fonttitle=\bfseries,
  title=#1
}

\newtcolorbox{info}[1][]{
  colback=blue!5!white,
  colframe=blue!75!black,
  fonttitle=\bfseries,
  title=#1
}

\newtcolorbox{astuce}[1][]{
  colback=green!5!white,
  colframe=green!75!black,
  fonttitle=\bfseries,
  title=#1
}

\newtcolorbox{exemple}[1][]{
  colback=orange!5!white,
  colframe=orange!75!black,
  fonttitle=\bfseries,
  title=#1
}

\newtcolorbox{schema}[1][]{
  colback=gray!5!white,
  colframe=gray!60!black,
  fonttitle=\bfseries,
  title=#1
}

\begin{document}

% ============================================================
% PAGE DE GARDE
% ============================================================
\begin{titlepage}
    \centering
    \vspace*{2cm}

    {\Huge \textbf{Documentation Technique — Phase 2}}

    \vspace{0.5cm}
    {\Large \textbf{ChatBot SUP'ONE}}

    \vspace{1.5cm}

    \begin{minipage}{0.8\textwidth}
        \centering
        \small Document de référence pour l'équipe projet \\
        expliquant les limites de la Phase 1, les nouveaux concepts \\
        introduits en Phase 2, le fonctionnement du pipeline RAG \\
        (Embeddings, FAISS, LLM) et le passage à React \\
    \end{minipage}

\end{titlepage}

% Table des matières
\tableofcontents
\clearpage

% ============================================================
\section{Introduction}
% ============================================================

\subsection{Pourquoi ce document existe}

Le Club Informatique SUP'PTIC a livré une première version fonctionnelle
du chatbot en Phase 1. Ce premier système répondait aux questions des
étudiants en cherchant dans une base de données la question la plus
proche de celle posée. Il a bien fonctionné --- mais l'utilisation réelle
a mis en évidence des limites concrètes qu'on ne pouvait pas contourner
sans changer d'approche.

Ce document explique honnêtement ce qui n'allait pas, pourquoi on a
choisi une nouvelle architecture, et comment elle fonctionne. Il n'y a
pas besoin d'être développeur pour le lire --- chaque concept technique
est expliqué avec des mots simples et des exemples tirés du projet.

\subsection{Ce que vous allez comprendre}

À la fin de ce document, vous saurez répondre à ces questions :

\begin{enumerate}
    \item Pourquoi la Phase 1 avait des limites qu'on ne pouvait pas
    corriger juste en ajoutant plus de données ?
    \item C'est quoi concrètement un embedding, FAISS, un LLM, le RAG ?
    \item Comment une question posée par un étudiant devient une réponse
    affichée à l'écran, étape par étape ?
    \item Qu'est-ce qui change dans le projet par rapport à ce qu'on avait ?
\end{enumerate}

\subsection{Public cible}

Ce document s'adresse à tous les membres du club --- développeurs,
data curators, designers, nouveaux arrivants. Aucune connaissance
préalable en intelligence artificielle n'est requise.

\clearpage

% ============================================================
\section{Ce qu'on avait en Phase 1}
% ============================================================

\subsection{L'architecture en quelques mots}

La Phase 1 était construite autour de \textbf{Django} et d'une base
de données \textbf{SQLite}. Les données --- questions, réponses,
catégories --- étaient stockées dans des fichiers JSON, puis importées
dans la base via un script Python. Une fois en base, Django s'occupait
de tout : stocker les FAQ, calculer les vecteurs de recherche,
répondre aux requêtes de l'interface.

L'algorithme de recherche utilisé était le \textbf{TF-IDF} combiné
à la \textbf{similarité cosinus}. En résumé : chaque question de la
base était transformée en une suite de nombres (un vecteur) représentant
les mots qu'elle contient. Quand un étudiant posait une question,
cette question était elle aussi transformée en vecteur, et le système
cherchait dans la base le vecteur le plus proche. Il retournait ensuite
la réponse associée.

\begin{info}[Ce que faisait le système à chaque question]
\begin{enumerate}
    \item L'étudiant envoie sa question via l'interface web
    \item Django transforme la question en vecteur TF-IDF
    \item Le système compare ce vecteur à tous ceux stockés en base
    \item Il trouve la question la plus proche et retourne sa réponse
    \item La réponse est affichée --- telle quelle, sans modification
\end{enumerate}
\end{info}

\subsection{Ce qui fonctionnait bien}

Le système Phase 1 avait de vraies qualités. Il était rapide,
stable, déployé et utilisable. L'API Django répondait correctement,
l'interface web était fonctionnelle, et le système de feedback
permettait de collecter les retours des utilisateurs. Pour des
questions formulées exactement comme celles stockées en base,
les résultats étaient très bons.

\subsection{Ce qui posait problème}

Le problème fondamental de TF-IDF, c'est qu'il compare des
\textbf{mots}, pas des \textbf{idées}. Deux phrases peuvent vouloir
dire la même chose avec des mots complètement différents --- et
TF-IDF les traitera comme étrangères l'une à l'autre.

\begin{exemple}[Le problème central de la Phase 1]
Un étudiant tape : \textit{«~C'est combien pour s'inscrire ?~»} \\
La base contient : \textit{«~Quels sont les frais de scolarité ?~»} \\
\vspace{0.3cm}
Ces deux phrases veulent dire la même chose. Mais elles n'ont
\textbf{aucun mot en commun}. Le score TF-IDF est donc proche
de zéro, et le chatbot répond à côté --- ou pire, dit qu'il ne
sait pas --- alors que la réponse est bien là, dans la base.
\end{exemple}

Ce problème ne pouvait pas être résolu en ajoutant plus de données.
Même avec 10 000 questions en base, si l'étudiant utilisait des
mots différents de ceux stockés, le système échouait. C'est une
limite \textbf{structurelle} de l'algorithme, pas un manque de contenu.

\clearpage

% ============================================================
\section{Les limites qui ont motivé le passage à la Phase 2}
% ============================================================

Voici les cinq limites concrètes observées après le déploiement
de la Phase 1.

\subsection{Limite 1 --- Le système comparait des mots, pas du sens}

C'est la limite la plus importante. TF-IDF ne comprend pas le sens
des phrases. Il compte les mots, calcule leur fréquence, et cherche
des correspondances lexicales. Deux synonymes, deux formulations
différentes d'une même idée, une faute de frappe --- et le résultat
devient mauvais. Un étudiant qui écrit \textit{«~tarif scolarité~»},
\textit{«~combien ça coûte~»} ou \textit{«~prix inscription~»} pour
demander la même chose obtiendra des résultats différents selon
la formulation, et pas toujours le bon.

\subsection{Limite 2 --- Les réponses étaient figées}

Le chatbot ne générait rien. Il copiait-collait directement le
contenu du champ \texttt{responses} de l'objet JSON le mieux
classé. Peu importe comment l'étudiant avait formulé sa question,
la réponse était toujours identique, rédigée d'une seule façon.
C'est fonctionnel, mais ça donne un résultat robotique --- et
surtout, ça ne s'adapte pas au contexte de la question posée.

\subsection{Limite 3 --- Impossible de croiser plusieurs sources}

Certaines questions touchent à plusieurs sujets à la fois.
\textit{«~Comment rejoindre le club informatique et quels sont
les frais ?~»} nécessite des informations qui viennent de deux
objets JSON différents. En Phase 1, le système retournait
toujours \textbf{une seule réponse} --- la meilleure correspondance
unique. Tout le reste était ignoré. Les réponses à des questions
un peu complexes étaient donc forcément incomplètes.

\subsection{Limite 4 --- Aucune mémoire entre les messages}

Chaque question était traitée de façon totalement indépendante.
Si un étudiant demandait d'abord \textit{«~Quelles sont les filières
disponibles ?~»}, puis enchaînait avec \textit{«~Et les frais pour
celle-là ?~»}, le chatbot ne savait pas de quelle filière il
parlait. Il repartait de zéro à chaque message, sans aucun
contexte de ce qui venait d'être échangé.

\subsection{Limite 5 --- Fragilité face aux variantes de langage}

Les abréviations, les fautes de frappe courantes, les expressions
familières ou les mots en franglais perturbaient le calcul TF-IDF.
\textit{«~suptptic~»}, \textit{«~frais scol~»}, \textit{«~comment
register~»} --- autant de formulations réelles que des étudiants
utilisent et qui faisaient échouer la recherche malgré la
présence de la réponse en base.

\vspace{0.5cm}

\begin{important}[Ce qu'il faut retenir]
Ces cinq limites ne venaient pas d'un manque de données
ni d'un bug dans le code. Elles venaient du fait que
\textbf{TF-IDF compare des mots, pas des intentions}.
Quelle que soit la quantité de données ajoutée, on ne
pouvait pas résoudre ce problème sans changer d'algorithme.
C'est pourquoi la Phase 2 introduit une approche différente.
\end{important}

\clearpage

% ============================================================
\section{Les concepts clés de la Phase 2}
% ============================================================

La Phase 2 introduit quatre notions nouvelles qui reviennent
constamment dans les discussions du projet. Il est important
de bien les distinguer car elles sont souvent confondues.

% ------------------------------------------------------------
\subsection{Les Embeddings --- lire le sens, pas les mots}
% ------------------------------------------------------------

En Phase 1, TF-IDF transformait une phrase en vecteur en
comptant les mots qu'elle contient. Deux phrases sans mot
en commun donnaient deux vecteurs très éloignés, même si
elles voulaient dire la même chose.

Les \textbf{embeddings} font quelque chose de fondamentalement
différent. Ils transforment une phrase en vecteur en se basant
sur son \textbf{sens}. Ils ont été produits par un modèle entraîné
sur des centaines de millions de phrases pour apprendre que
\textit{«~frais de scolarité~»} et \textit{«~combien coûte
l'inscription~»} expriment la même intention --- et donc produisent
des vecteurs proches.

\begin{exemple}[La même question avec les deux algorithmes]
\textbf{Question posée :} \textit{«~C'est combien pour s'inscrire ?~»} \\
\textbf{Question en base :} \textit{«~Quels sont les frais de scolarité ?~»} \\
\vspace{0.3cm}
Avec TF-IDF (Phase 1) : score $\approx$ 0.00 --- aucun mot en commun \\
Avec Embeddings (Phase 2) : score $\approx$ 0.82 --- même intention détectée
\end{exemple}

Le modèle utilisé s'appelle \textbf{MiniLM-L6-v2}. Ce n'est pas
du code que l'équipe écrit --- c'est une bibliothèque Python
(\texttt{pip install sentence-transformers}) que l'on installe
et que l'on utilise. Elle transforme n'importe quel texte en
un vecteur de \textbf{384 nombres}.

\begin{astuce}[Analogie simple]
Imaginez que chaque phrase reçoit des coordonnées GPS
dans un espace à 384 dimensions. Deux phrases qui veulent
dire la même chose se retrouvent proches géographiquement,
même si leurs mots sont différents. TF-IDF, lui, regardait
uniquement les mots --- pas la position dans cet espace.
\end{astuce}

% ------------------------------------------------------------
\subsection{FAISS --- chercher vite parmi des milliers de vecteurs}
% ------------------------------------------------------------

Une fois que toutes les questions de la base ont été transformées
en vecteurs par MiniLM, il faut pouvoir chercher très rapidement
lequel est le plus proche d'une nouvelle question. C'est exactement
ce que fait \textbf{FAISS}.

FAISS est une bibliothèque développée par Meta (Facebook),
installée via \texttt{pip install faiss-cpu}. Elle stocke
tous les vecteurs dans un fichier \texttt{index.bin} et
permet de retrouver en quelques millisecondes les cinq vecteurs
les plus proches d'un vecteur donné, même parmi des dizaines
de milliers.

FAISS ne sait pas ce que veulent dire les vecteurs. Il sait
juste les comparer très vite. C'est le rôle de \texttt{metadata.json}
de faire le lien entre un numéro de vecteur et la réponse
correspondante dans les fichiers JSON.

\begin{schema}[Ce que FAISS reçoit et ce qu'il retourne]
\textbf{Entrée :} le vecteur de la question posée par l'étudiant \\
\textbf{Sortie :} les numéros des 5 vecteurs les plus proches \\
\vspace{0.3cm}
Exemple : \texttt{[42, 871, 304, 1205, 2087]} \\
\vspace{0.3cm}
\texttt{metadata.json} est alors consulté pour retrouver
la réponse associée à chacun de ces numéros.
\end{schema}

% ------------------------------------------------------------
\subsection{Le LLM --- rédiger une réponse naturelle}
% ------------------------------------------------------------

En Phase 1, une fois la meilleure correspondance trouvée,
le système copiait-collait la réponse telle quelle depuis
la base. En Phase 2, on introduit un \textbf{LLM} (Large
Language Model) --- un modèle d'intelligence artificielle
capable de lire un texte et de rédiger une réponse naturelle.

Le LLM utilisé s'appelle \textbf{Phi-3 Mini}. Il tourne
entièrement sur la machine locale via un outil appelé
\textbf{Ollama} --- sans connexion internet, sans envoyer
de données à l'extérieur.

Concrètement, le LLM reçoit un message structuré contenant :
la question posée par l'étudiant, et les cinq réponses
récupérées par FAISS. Il rédige alors une réponse fluide
et naturelle à partir de ces informations.

\begin{important}[Le LLM ne devine rien --- il s'appuie sur vos données]
Le LLM ne puise pas dans ses propres connaissances générales.
Il ne peut répondre qu'avec ce qu'on lui donne. Si FAISS
ne trouve aucune correspondance pertinente, le système
répond : \textit{«~Je n'ai pas cette information ---
contactez le secrétariat.~»} \\
\vspace{0.2cm}
C'est une garantie importante : le chatbot ne peut pas
inventer des informations sur SUP'PTIC.
\end{important}

% ------------------------------------------------------------
\subsection{Le RAG --- le nom de la stratégie globale}
% ------------------------------------------------------------

\textbf{RAG} signifie \textit{Retrieval Augmented Generation}.
Ce n'est pas une bibliothèque, pas un outil à installer.
C'est le nom qu'on donne à la stratégie qui combine les trois
éléments précédents :

\begin{enumerate}
    \item \textbf{Retrieval} --- on \textit{récupère} les réponses
    les plus pertinentes grâce aux embeddings et à FAISS
    \item \textbf{Augmented} --- on \textit{enrichit} le contexte
    donné au LLM avec ces réponses
    \item \textbf{Generation} --- le LLM \textit{génère} une réponse
    naturelle à partir de ce contexte
\end{enumerate}

Le RAG est la recette. MiniLM, FAISS et Phi-3 sont les ingrédients.
Le code Django de l'équipe est le cuisinier qui suit la recette.

\clearpage

% ============================================================
\section{Le flux complet --- de la question à la réponse}
% ============================================================

\subsection{Vue d'ensemble}

Le système Phase 2 fonctionne en deux temps bien distincts.
Le premier est un travail préparatoire fait une seule fois
par l'équipe technique. Le second se répète automatiquement
à chaque message d'un étudiant.

\begin{schema}[Les deux temps du système]
\textbf{Temps 1 --- Préparation (une seule fois)} \\
\hspace{1cm} Les fichiers JSON sont lus, vectorisés et indexés. \\
\hspace{1cm} Le système est prêt à recevoir des questions. \\
\vspace{0.4cm}
\textbf{Temps 2 --- Réponse (à chaque message)} \\
\hspace{1cm} La question est vectorisée, les réponses proches sont
trouvées, \\
\hspace{1cm} le LLM rédige et l'interface affiche.
\end{schema}

% ------------------------------------------------------------
\subsection{Temps 1 --- La préparation de l'index}
% ------------------------------------------------------------

Cette étape est invisible pour l'utilisateur final. Elle est
exécutée par l'équipe technique chaque fois que de nouveaux
contenus sont ajoutés aux fichiers JSON.

\vspace{0.3cm}
\textbf{Étape 1 --- Lecture des fichiers JSON}

Le script Python parcourt tous les fichiers JSON du dossier
\texttt{data/}. Pour chaque objet JSON, il extrait tous les
éléments de la liste \texttt{examples} --- ce sont les différentes
façons de formuler une question sur un même sujet.

\vspace{0.3cm}
\textbf{Étape 2 --- Vectorisation par MiniLM}

Chaque example est envoyé à MiniLM qui le transforme en un
vecteur de 384 nombres. Un objet JSON avec 8 examples
produit donc 8 vecteurs distincts --- mais tous les 8 pointent
vers la même réponse, puisqu'ils viennent du même objet.

\vspace{0.3cm}
\textbf{Étape 3 --- Construction de l'index FAISS}

Tous les vecteurs produits sont stockés dans \texttt{index.bin}.
En parallèle, \texttt{metadata.json} est écrit : il associe
chaque numéro de vecteur à la réponse et aux informations
(catégorie, source) de l'objet JSON dont il est issu.

\begin{exemple}[Ce que contient metadata.json pour 3 vecteurs du même objet]
\begin{lstlisting}[language=]
{ "vecteur_id": 0,
  "intent": "faq_frais_0",
  "example": "C'est combien pour s'inscrire ?",
  "response": "Les frais sont de 500 000 FCFA par an...",
  "categorie": "Admissions", "source": "e-supptic.cm" },

{ "vecteur_id": 1,
  "intent": "faq_frais_0",
  "example": "Quels sont les frais de scolarite ?",
  "response": "Les frais sont de 500 000 FCFA par an...",
  "categorie": "Admissions", "source": "e-supptic.cm" },

{ "vecteur_id": 2,
  "intent": "faq_frais_0",
  "example": "Tarif inscription SUP'PTIC",
  "response": "Les frais sont de 500 000 FCFA par an...",
  "categorie": "Admissions", "source": "e-supptic.cm" }
\end{lstlisting}
Les vecteurs 0, 1 et 2 ont des examples différents mais
la même réponse --- ils viennent du même objet JSON.
\end{exemple}

% ------------------------------------------------------------
\subsection{Temps 2 --- Le traitement de chaque question}
% ------------------------------------------------------------

\begin{exemple}[Un étudiant tape : «~C'est combien pour s'inscrire ?~»]

\textbf{Étape 1 --- Réception par Django} \\
L'interface envoie la question à l'API Django via une requête HTTP.

\vspace{0.4cm}
\textbf{Étape 2 --- Vectorisation de la question} \\
MiniLM transforme la question en un vecteur de 384 nombres,
exactement comme il l'a fait pour les examples lors de la
préparation.

\vspace{0.4cm}
\textbf{Étape 3 --- Recherche dans FAISS} \\
FAISS compare ce vecteur à tous ceux stockés dans
\texttt{index.bin} et retourne les 5 numéros les plus proches.

\vspace{0.4cm}
\textbf{Étape 4 --- Vérification de la pertinence} \\
Chaque résultat a un score entre 0 et 1. Trois cas possibles :
\begin{itemize}
    \item Score $<$ 0.30 : aucun résultat RAG pertinent ---
    \textbf{le système bascule automatiquement sur le moteur
    TF-IDF de la Phase 1} qui répond comme avant.
    L'étudiant obtient toujours une réponse, sans voir
    la différence.
    \item Score entre 0.30 et 0.55 : résultat approximatif ---
    la réponse sera précédée d'un avertissement
    \item Score $>$ 0.55 : bonne correspondance --- le LLM rédige
    normalement
\end{itemize}

\vspace{0.4cm}
\textbf{Étape 5 --- Construction du message pour le LLM} \\
Le système prépare un message structuré qui contient :
le rôle du chatbot, les 5 réponses récupérées par FAISS
avec leurs métadonnées, et la question originale de l'étudiant.
La dernière instruction est toujours : \textit{«~Réponds
uniquement à partir des informations fournies.~»}

\vspace{0.4cm}
\textbf{Étape 6 --- Génération par Phi-3 Mini} \\
Phi-3 lit le message et rédige une réponse mot par mot.
La réponse est envoyée en \textbf{streaming} : chaque mot
est transmis à l'interface dès qu'il est produit, ce qui
donne l'impression que le chatbot tape en temps réel.

\vspace{0.4cm}
\textbf{Étape 7 --- Affichage} \\
L'interface affiche la réponse progressivement. En dessous,
les 3 meilleures sources sont affichées avec leur catégorie
et leur score, pour que l'étudiant sache d'où vient l'information.

\vspace{0.4cm}
\textbf{Résultat affiché :} \\
\textit{«~Les frais de scolarité à SUP'PTIC s'élèvent à 500~000~FCFA
par an. Ce montant couvre l'ensemble de la formation. Des facilités
de paiement peuvent être demandées à la scolarité.~»} \\
\vspace{0.1cm}
\textit{Sources : Admissions (0.91) | e-supptic.cm}

\end{exemple}

\clearpage

\subsection{Schéma du flux complet}

\begin{figure}[H]
\centering
\begin{tikzpicture}[
    scale=0.78,
    every node/.style={transform shape},
    node distance=0.75cm,
    box/.style={rectangle, draw, thick, rounded corners=4pt,
                minimum width=9cm, minimum height=0.9cm,
                align=center, font=\small},
    offline/.style={box, fill=blue!8},
    online/.style={box, fill=orange!8},
    decision/.style={box, fill=red!8},
    result/.style={box, fill=green!8},
    arrow/.style={->, thick, >=stealth}
]

% Phase Offline
\node[offline] (json)
    {Lecture des fichiers JSON --- extraction des \texttt{examples}};
\node[offline, below=of json] (embed_off)
    {MiniLM --- chaque example devient un vecteur de 384 nombres};
\node[offline, below=of embed_off] (faiss_build)
    {FAISS construit \texttt{index.bin} \quad + \quad \texttt{metadata.json} est écrit};

% Séparateur
\node[below=0.5cm of faiss_build,
      font=\small\itshape, text=gray!70] (sep)
    {--- Système prêt --- se répète à chaque question ---};

% Phase Online
\node[online, below=0.5cm of sep] (question)
    {Question de l'étudiant reçue par l'API Django};
\node[online, below=of question] (embed_on)
    {MiniLM --- la question devient un vecteur de 384 nombres};
\node[online, below=of embed_on] (search)
    {FAISS cherche dans \texttt{index.bin} --- retourne le top-5};
\node[decision, below=of search] (seuil)
    {Vérification du score --- pertinent ou pas ?};
\node[online, below=of seuil] (prompt)
    {Construction du message --- question + 5 réponses + métadonnées};
\node[online, below=of prompt] (llm)
    {Phi-3 Mini rédige une réponse naturelle (streaming)};
\node[result, below=of llm] (affichage)
    {Affichage dans l'interface --- réponse + sources};

% Nœud fallback TF-IDF (à droite du nœud seuil)
\node[box, fill=purple!8, right=1.5cm of seuil] (fallback)
    {Fallback TF-IDF (Phase 1) ---\\réponse directe depuis la base};

% Flèches Offline
\draw[arrow] (json) -- (embed_off);
\draw[arrow] (embed_off) -- (faiss_build);

% Flèches Online
\draw[arrow] (question) -- (embed_on);
\draw[arrow] (embed_on) -- (search);
\draw[arrow] (search) -- (seuil);
\draw[arrow] (seuil) -- (prompt);
\draw[arrow] (prompt) -- (llm);
\draw[arrow] (llm) -- (affichage);

% Flèche fallback avec label
\draw[arrow, dashed, color=purple!70!black]
    (seuil) -- (fallback)
    node[midway, above, font=\tiny\itshape, text=purple!70!black]
    {score $<$ 0.30};

% Grouper les boîtes Offline
\node[draw=none, fit=(json)(faiss_build)] (offlineGroup) {};
\draw[thick, decorate,
      decoration={brace, amplitude=7pt},
      xshift=-1cm] % petit décalage vers la gauche
    (offlineGroup.south west) -- (offlineGroup.north west)
    node[midway, left=10pt,
         font=\small\bfseries, text=blue!60!black]
    {Préparation};

% Grouper les boîtes Online
\node[draw=none, fit=(question)(affichage)] (onlineGroup) {};
\draw[thick, decorate,
      decoration={brace, amplitude=7pt},
      xshift=-1cm]
    (onlineGroup.south west) -- (onlineGroup.north west)
    node[midway, left=10pt,
         font=\small\bfseries, text=orange!70!black]
    {À chaque question};


\end{tikzpicture}
\caption{Flux complet du système RAG --- Phase 2}
\end{figure}

\clearpage

% ============================================================
\section{Ce qui change dans le projet}
% ============================================================

\subsection{Django ne disparaît pas --- son rôle se réduit}

En Phase 1, Django était le cœur du système : il stockait
les FAQ, calculait les vecteurs TF-IDF, et faisait la recherche.
Tout passait par la base SQLite.

En Phase 2, la recherche et la génération de réponses sont
prises en charge par FAISS et Phi-3. Django conserve ce
qu'il fait naturellement bien : gérer les utilisateurs,
enregistrer les feedbacks, exposer les endpoints de l'API,
et afficher les statistiques.

\begin{table}[H]
\centering
\begin{tabularx}{\textwidth}{|l|X|X|}
\hline
\textbf{Responsabilité} & \textbf{Phase 1} & \textbf{Phase 2} \\
\hline
Stocker les FAQ & Base Django/SQLite & Fichiers JSON sur le disque \\
\hline
Stocker les vecteurs & Table \texttt{FAQVector} en base & \texttt{index.bin} (FAISS) \\
\hline
Faire la recherche & TF-IDF sur la base & MiniLM + FAISS \\
\hline
Produire la réponse & Copier depuis la base & Phi-3 Mini génère \\
\hline
\textbf{Fallback si RAG insuffisant} & \textit{n/a} &
\textbf{TF-IDF Phase 1 --- bascule automatique si score FAISS $<$ 0.30} \\
\hline
Feedbacks & Base Django & Base Django (inchangé) \\
\hline
Utilisateurs & Base Django & Base Django (inchangé) \\
\hline
Statistiques & Base Django & Base Django (inchangé) \\
\hline
API REST & Django REST Framework & Django REST Framework (inchangé) \\
\hline
\end{tabularx}
\caption{Répartition des responsabilités --- Phase 1 vs Phase 2}
\end{table}

\subsection{Les nouveaux fichiers qui entrent en jeu}

Deux fichiers nouveaux font leur apparition dans le projet :

\begin{itemize}
    \item \texttt{index.bin} --- c'est l'index FAISS. Il contient
    tous les vecteurs générés à partir des examples JSON.
    C'est lui que FAISS consulte lors de chaque recherche.
    Il est regénéré chaque fois qu'on ajoute de nouveaux
    contenus aux fichiers JSON.

    \item \texttt{metadata.json} --- c'est l'annuaire qui fait
    le lien entre un numéro de vecteur dans FAISS et la réponse
    correspondante dans les fichiers JSON. Sans lui, FAISS
    retournerait des numéros sans pouvoir les associer à du contenu.
\end{itemize}

\subsection{L'endpoint principal change de comportement}

L'endpoint \texttt{POST /api/chatbot/ask/} existait déjà en
Phase 1. En Phase 2, son comportement interne change complètement :
au lieu d'interroger la base SQLite avec TF-IDF, il lance le
pipeline RAG --- vectorisation de la question, recherche FAISS,
construction du message, appel à Ollama, streaming de la réponse.

Deux nouveaux endpoints font leur apparition :

\begin{itemize}
    \item \texttt{GET /api/sources/} --- retourne les sources JSON
    utilisées pour produire la dernière réponse, pour affichage
    dans l'interface
    \item \texttt{POST /api/reload-index/} --- recharge \texttt{index.bin}
    en mémoire sans redémarrer Django, utile après l'ajout de
    nouveaux contenus JSON
\end{itemize}

\clearpage

% ============================================================
\section{Migration de l'interface : de la PWA à React}
% ============================================================

\subsection{Ce qu'était la PWA en Phase 1}

En Phase 1, l'interface utilisateur était une \textbf{Progressive Web App}
(PWA) développée en HTML, CSS et JavaScript vanilla. Elle était servie
par un petit serveur Python indépendant (\texttt{start\_server.py}) sur
le port 8080, séparé du backend Django. Un \texttt{service-worker.js}
lui permettait de fonctionner partiellement hors-ligne et d'être
installable sur mobile.

Cette approche était simple à mettre en place et fonctionnelle pour
une première version. Mais elle montrait ses limites dès que
l'interface devenait plus complexe : gestion d'état dispersée dans
le DOM, code JavaScript difficile à structurer au fil des nouvelles
fonctionnalités, réutilisation des composants quasi impossible.

\subsection{Pourquoi passer à React}

La Phase 2 introduit plusieurs nouveautés côté interface qui auraient
été très difficiles à gérer proprement en JavaScript vanilla :

\begin{itemize}
    \item \textbf{Le streaming de réponse} --- afficher les mots du LLM
    au fur et à mesure qu'ils arrivent nécessite une gestion d'état
    réactive. En React, un simple \texttt{useState} suffit ; en JS
    vanilla, cela demande une manipulation du DOM fastidieuse et
    fragile.

    \item \textbf{L'affichage des sources} --- chaque réponse est
    accompagnée de 3 sources avec catégorie et score. Ce sont des
    composants répétables, exactement le cas d'usage pour lequel
    React a été conçu.

    \item \textbf{La gestion de l'historique de conversation} ---
    maintenir et afficher une liste de messages avec états différents
    (en cours, répondu, erreur) est naturel en React grâce à la
    notion de composants et de props.

    \item \textbf{L'indicateur de confiance} --- afficher visuellement
    le niveau de pertinence d'une réponse (score FAISS) sous forme
    de jauge ou de badge coloré est trivial en JSX, laborieux en
    HTML/JS manuel.
\end{itemize}

\begin{info}[Ce que React change fondamentalement]
En PWA vanilla, l'interface \textit{manipule} le DOM pour refléter
l'état de l'application. En React, l'interface \textit{est} une
fonction de l'état --- on met à jour l'état, React s'occupe
du rendu. Pour une interface de chat avec des données dynamiques
qui arrivent en streaming, c'est un gain de fiabilité et de
lisibilité considérable.
\end{info}

\subsection{Ce qui disparaît et ce qui reste}

La migration ne repart pas de zéro. La logique métier reste
intégralement côté Django --- React ne remplace que la couche
d'affichage.

\begin{table}[H]
\centering
\begin{tabularx}{\textwidth}{|l|X|X|}
\hline
\textbf{Élément} & \textbf{Phase 1 (PWA)} & \textbf{Phase 2 (React)} \\
\hline
Technologie & HTML/CSS/JS vanilla & React (JSX + hooks) \\
\hline
Serveur frontend & \texttt{start\_server.py} sur port 8080 & \texttt{npm start} / Vite sur port 3000 \\
\hline
Gestion d'état & Variables JS + manipulation DOM & \texttt{useState}, \texttt{useEffect} \\
\hline
Composants & Fonctions JS manuelles & Composants React réutilisables \\
\hline
Streaming LLM & EventSource / fetch manuel & Hook dédié avec mise à jour réactive \\
\hline
Offline & Service Worker + cache & Non prioritaire en Phase 2 \\
\hline
Installation mobile & Manifest PWA & Non prioritaire en Phase 2 \\
\hline
Communication API & \texttt{fetch()} vers Django & \texttt{fetch()} vers Django (inchangé) \\
\hline
\end{tabularx}
\caption{Comparaison interface --- Phase 1 (PWA) vs Phase 2 (React)}
\end{table}

\begin{astuce}[Ce que ça change pour les autres équipes]
Pour l'équipe backend et data, \textbf{rien ne change} : l'API
Django REST reste identique, les endpoints sont les mêmes. Seule
l'équipe frontend travaille différemment --- elle passe de fichiers
\texttt{.html}/\texttt{.js} à des composants \texttt{.jsx} organisés
dans un projet React. Les fonctionnalités PWA (mode offline,
installation sur mobile) pourront être réintégrées ultérieurement
via des librairies React dédiées si le besoin s'en fait sentir.
\end{astuce}

\clearpage

% ============================================================
\section{Points clés à retenir}
% ============================================================

\begin{important}[Ce qu'il faut avoir compris à la fin de ce document]
\begin{enumerate}
    \item \textbf{La Phase 1 avait une limite structurelle.}
    TF-IDF compare des mots, pas des intentions. Ajouter plus
    de données ne pouvait pas résoudre ce problème.

    \item \textbf{Les embeddings remplacent TF-IDF.}
    MiniLM transforme les textes en vecteurs basés sur le sens.
    Deux phrases similaires en intention donnent des vecteurs
    proches, même sans mot en commun.

    \item \textbf{FAISS stocke et cherche dans ces vecteurs.}
    Ce n'est pas du code produit par l'équipe --- c'est une
    bibliothèque de Meta qu'on installe et qu'on utilise.
    Elle cherche en quelques millisecondes parmi des milliers
    de vecteurs.

    \item \textbf{MiniLM vectorise chaque example séparément.}
    Un objet JSON avec 8 examples produit 8 vecteurs dans FAISS,
    tous liés à la même réponse. Plus il y a d'examples,
    plus la recherche est précise.

    \item \textbf{metadata.json fait le lien entre les vecteurs
    et les réponses.} FAISS retourne des numéros --- metadata.json
    dit ce que chaque numéro représente.

    \item \textbf{Le LLM rédige, il n'invente pas.}
    Phi-3 Mini reçoit uniquement les réponses trouvées par FAISS.
    Il ne peut pas produire d'informations qui ne sont pas
    dans vos fichiers JSON.

    \item \textbf{RAG est le nom de la stratégie, pas un outil.}
    Chercher (FAISS) $\rightarrow$ enrichir (JSON)
    $\rightarrow$ générer (LLM) : voilà ce qu'est le RAG.

    \item \textbf{Django reste dans le projet.}
    Il garde les feedbacks, les utilisateurs, les statistiques
    et l'API REST. Ce qui change, c'est uniquement le moteur
    de recherche et la génération de réponse.

    \item \textbf{L'interface PWA est remplacée par React.}
    La PWA HTML/JS vanilla de la Phase 1 laisse place à une
    application React. Le streaming LLM, les sources et
    l'historique de conversation sont bien plus faciles à
    gérer avec des composants et des hooks. L'API Django
    ne change pas --- seule la couche d'affichage évolue.

    \item \textbf{TF-IDF n'est pas supprimé --- il devient le filet de sécurité.}
    Si le score FAISS est trop bas (question hors-domaine,
    formulation très inhabituelle), le système bascule
    automatiquement sur le moteur TF-IDF de la Phase 1.
    L'étudiant obtient toujours une réponse.
    \textbf{La Phase 2 ajoute de l'intelligence sans rien retirer.}
\end{enumerate}
\end{important}

    \vfill

    \begin{minipage}{0.8\textwidth}
        \centering
        \small
        \textit{Rédigé par le Chef de Division Programmation} \\
        \textit{Club Informatique SUP'PTIC --- Projet SUP'ONE}
    \end{minipage}

\end{document}